%---------------------------------------------------------------------
%	Abstract
%---------------------------------------------------------------------
\chapter{Abstract}				
\label{ch:abs}
%---------------------------------------------------------------------
% TAKEN FROM UCT WEBSITE
% The Doctoral Degrees Board recommends that candidates include an abstract 
% that fits onto one page, and which includes the author's full name, thesis 
% title and date. The text should not exceed 350 words. Candidates may also 
% include a more substantial summary of their work, in addition to the abstract.
% Follow the same format for an MSc. abstract.

% The abstract should stand on its own. It should answer the following questions:
%	1.	What did the author do? What ideas, notions, hypotheses, 
%		concepts, theories or thoughts were investigated?
%	2.	How did the author do the work? What data were generated and used? 
%		What was the origin of the data? How were data gathered? 
%		What tests, scales, indices, or summary measures were used? 
%		In other words, how was the analysis and/or synthesis done?
%	3.	What were the conclusions and what were the significant findings?
%
% Some studies cannot readily be summarised in this way and require more 
% descriptive abstracts. Do not use telegraphic phrases. Do not repeat 
% information given in the title. Do not use abbreviations. The purpose of 
% an abstract is to enable a researcher/examiner to understand the essential 
% hypothesis, method and findings of the research.
%---------------------------------------------------------------------
\begin{center}
	\textbf{\Large \titl}\\
			\vskip 0.2cm
			\auth\\
			\vskip 0.2cm
	\textit{\footnotesize \FullDate}
			\vskip 1cm
\end{center}
%---------------------------------------------------------------------
% Write abstract here [350 words]
%---------------------------------------------------------------------

%Sediment microtopography, and rugosity in particular, has been proposed as a non-destructive proxy for the bioturbation activity of the sand prawn Callichirus kraussi on the intertidal flats of Langebaan Lagoon. An existing scanning system based on an Intel RealSense D456 active stereo camera has been implemented, however has never been characterised for these environmental conditions, and it analysed each of its repeated multi-view scans individually. This project characterised the D456 for close-range scanning of sediment-like surfaces and developed a processing pipeline that crops, registers and fuses multi-view scans before computing ISO 25178-2 areal roughness parameters.
%
%Planar tests showed that noise increased with working distance, tilt and ambient light. Plane-fit RMS rose from 1.22\,mm at 604\,mm to 4.58\,mm at 903\,mm. Scans of 3D-printed reference objects showed that viewing geometry matters more than software settings. Increasing resolution to 1280×720 reduced RMSE by 13.4\%, and the High Density preset gave the best F-score. Post-processing filters improved RMSE by at most 0.24\,mm.
%
%The pipeline uses AprilTags and IMU gravity readings for coarse alignment, refines the alignment with point-to-plane ICP, and reconciles the overlapping scans with a confidence-weighted trimmed mean. Fusion improved every metric at every height. On the large reference surface, it raised the F-score by 14.0 percentage points, reaching 90.9\% at 641\,mm. Twelve well-spread views matched the accuracy of a full 44-view sweep.
%
%On flat laboratory sand, fused scans measured a rugosity of 1.0034, against 1.0193 for single scans. On shaped sand, single scans added 2.1 to 2.6 percentage points of spurious excess area, an amount comparable to the signal itself. Rugosity is therefore only comparable between surveys that use a fixed capture and processing protocol, and this report recommends such a protocol.

Sediment microtopography, and rugosity in particular, has been proposed as a non-destructive proxy  candidate for bioturbation by the sand prawn \textit{Kraussillichirus kraussi} on the intertidal flats of Langebaan Lagoon. An existing scanning system captures repeated multi-view scans of the sediment using an Intel RealSense D456 active stereo camera. However, the camera had not been characterised for close-range scanning of sediment, and each scan was analysed individually, so the reliability of the resulting rugosity values was unknown.

This project characterised the D456 using planar targets at varying working distance, tilt and ambient light, and using 3D-printed reference objects scanned at different resolutions, depth presets and post-processing filter settings. A pipeline was then developed that crops, registers and fuses the multi-view scans before computing ISO 25178-2 areal roughness parameters. AprilTag fiducials and gravity readings from the camera's inertial measurement unit provide coarse alignment. The Point-to-plane iterative closest point method refines this alignment, and a confidence-weighted trimmed mean reconciles the overlapping scans. The pipeline was evaluated on reference objects and on flat and shaped laboratory sand.

Noise increased with working distance, tilt and ambient light. The plane-fit root-mean-square residual rose from 1.22\,mm at 604\,mm to 4.58\,mm at 903\,mm. Viewing geometry mattered more than software settings. Increasing the resolution to 1280$\times$720 reduced error by 13.4\%, the High Density preset gave the best F-score, and post-processing filters improved error by at most 0.24\,mm. Fusion improved every metric at every height. On the large reference surface, it raised the F-score by 14.0 percentage points to 90.9\% at 641\,mm, and twelve well-spread views matched the accuracy of a full 44-view sweep. On flat sand, fused scans measured a rugosity of 1.0034, against 1.0193 for single scans. On shaped sand, single scans added 2.1 to 2.6 percentage points of spurious excess area, an amount comparable to the signal itself.

Single-scan rugosity is therefore inflated by sensor noise, and multi-view fusion substantially reduces this bias. Rugosity values are only comparable between surveys that share a fixed capture and processing protocol. It is recommended that surveys scan at the shortest practical working distance, use 1280$\times$720 resolution with the High Density preset, capture at least twelve well-spread views, and report rugosity from fused rather than single scans.
